%% =============================================================================
%% sm_c5_cholesky.tex — SM-C, Subsection C.5: Cholesky Affine Transformation
%% Body content only (\subsection level); parent structure in sm_c.tex
%% Primary source: dev/log/06_research-notes-survey-design-pseudo-posterior.qmd
%%   lines 1197--1261 (theorem proof, block-wise application, limitations)
%% Corrections applied: T1-4 (MCMC vs design-consistent shrinkage),
%%   T2-1 (block-wise Cholesky caveat)
%% =============================================================================

\subsection{Cholesky affine transformation}
\label{smc:cholesky}

The sandwich variance $\mathbf{V}_{\mathrm{sand}}$~\eqref{eq:V-sand}
provides the design-consistent covariance, but the MCMC draws
$\{\btheta^{(m)}\}_{m=1}^M$ from the pseudo-posterior have empirical
covariance $\bSigma_{\mathrm{MCMC}} \approx \mathbf{H}_{\mathrm{obs}}^{-1}$,
not $\mathbf{V}_{\mathrm{sand}}$.  In a frequentist setting one would simply
report $\mathbf{V}_{\mathrm{sand}}$ as the variance estimator; in the
Bayesian framework, we require the posterior draws themselves to have the
correct covariance so that credible intervals are properly calibrated
\citep{WilliamsSavitsky2021}.
\cref{thm:cholesky} in the main text states the affine transformation
that achieves this.  We now provide the full proof.

%% ---------------------------------------------------------------------------
%%  Proof of Theorem (thm:cholesky)
%% ---------------------------------------------------------------------------

\begin{proof}[Proof of \textup{\cref{thm:cholesky}}]
  Define centered draws
  $\mathbf{u}^{(m)} = \btheta^{(m)} - \hat{\btheta}$, so that
  $\bSigma_{\mathrm{MCMC}}
  = M^{-1}\sum_{m=1}^{M}\mathbf{u}^{(m)}(\mathbf{u}^{(m)})\t$.
  Set
  \begin{equation}\label{smc:eq-A}
    \mathbf{A}
    = \mathbf{L}_{\mathrm{sand}}\,\mathbf{L}_{\mathrm{MCMC}}^{-1},
  \end{equation}
  where $\mathbf{L}_{\mathrm{MCMC}}
  = \mathrm{chol}(\bSigma_{\mathrm{MCMC}})$ and
  $\mathbf{L}_{\mathrm{sand}}
  = \mathrm{chol}(\mathbf{V}_{\mathrm{sand}})$ are lower-triangular
  Cholesky factors.  The transformation~\eqref{eq:cholesky-transform} is
  then $\btheta^{*(m)} = \hat{\btheta} + \mathbf{A}\,\mathbf{u}^{(m)}$.

  \medskip\noindent
  \textup{(a) Mean preservation.}
  \[
    \frac{1}{M}\sum_{m=1}^{M}\btheta^{*(m)}
    = \hat{\btheta}
      + \mathbf{A}\,\Bigl(\frac{1}{M}\sum_{m=1}^{M}\mathbf{u}^{(m)}\Bigr)
    = \hat{\btheta} + \mathbf{A}\cdot\mathbf{0}
    = \hat{\btheta}.
  \]

  \noindent
  \textup{(b) Covariance recovery.}
  \begin{align}
    \frac{1}{M}\sum_{m=1}^{M}
      (\btheta^{*(m)} - \hat{\btheta})\,
      (\btheta^{*(m)} - \hat{\btheta})\t
    &= \mathbf{A}\,\bSigma_{\mathrm{MCMC}}\,\mathbf{A}\t
      \notag \\
    &= \mathbf{L}_{\mathrm{sand}}\,
       \mathbf{L}_{\mathrm{MCMC}}^{-1}\,
       \mathbf{L}_{\mathrm{MCMC}}\,
       \mathbf{L}_{\mathrm{MCMC}}\t\,
       \mathbf{L}_{\mathrm{MCMC}}^{-\t}\,
       \mathbf{L}_{\mathrm{sand}}\t
      \notag \\
    &= \mathbf{L}_{\mathrm{sand}}\,\mathbf{I}\,
       \mathbf{L}_{\mathrm{sand}}\t
     = \mathbf{V}_{\mathrm{sand}}.
      \label{smc:eq-cholesky-var}
  \end{align}

  \noindent
  \textup{(c) Affine equivariance.}
  The standardized draws
  $\mathbf{z}^{(m)} = \mathbf{L}_{\mathrm{MCMC}}^{-1}\mathbf{u}^{(m)}$
  satisfy
  \[
    M^{-1}\sum_m \mathbf{z}^{(m)}(\mathbf{z}^{(m)})\t = \mathbf{I}.
  \]
  The transformation maps
  \[
    \btheta^{*(m)} - \hat{\btheta}
    = \mathbf{L}_{\mathrm{sand}}\,\mathbf{z}^{(m)},
  \]
  which is a linear rescaling of the standardized draws.  The
  Wald confidence intervals---computed as
  $\hat{\theta}_p \pm z_{\alpha/2}\sqrt{[\mathbf{V}_{\mathrm{sand}}]_{pp}}$---depend
  only on the marginal variances of~$\btheta^*$ and are therefore
  invariant to the choice of~$\mathbf{A}$.  Marginal quantiles of
  the transformed draws may differ from those of the original
  draws when~$\mathbf{A}$ is non-diagonal.
\end{proof}

%% ---------------------------------------------------------------------------
%%  Block-wise application
%% ---------------------------------------------------------------------------

\begin{remark}[Block-wise application]
\label{smc:rem-blockwise}
  In practice the transformation~\eqref{eq:cholesky-transform} is applied
  block-wise rather than to the full
  $\approx\!616$-dimensional vector simultaneously.  The three
  parameter classes are listed below; the numerical Cholesky procedure
  operates on blocks~(i) and~(ii) only, since block~(iii) receives no
  correction
  (\cref{rem:blockwise} in the main text).  The rationale follows from
  \cref{smc:prop-H-blockdiag} and~\cref{smc:rem-J-nonblockdiag}\textup{:}
  \begin{enumerate}[label=\textup{(\roman*)},nosep]
    \item \emph{Fixed effects}
      $\boldsymbol{\xi}
      = (\balpha\t,\bbeta\t,\log\kappa)\t \in \R^{11}$\textup{:}
      full Cholesky correction.  The $\DER$ ranges from $1.14$ to
      $4.18$ (mean $2.11$), and the corresponding confidence-interval
      inflation factors are $\sqrt{\DER} \in [1.07, 2.04]$.
    \item \emph{State random effects}
      $\bepsilon_s \in \R^{10}$, $s = 1,\ldots,S$\textup{:} partial
      correction.  The posterior variance of $\bepsilon_s$ mixes the
      likelihood information from providers in state~$s$ (affected by
      survey weights) with the shrinkage toward the prior mean
      $\bGamma_k\mathbf{v}_s$ (unaffected).  When the prior
      dominates (small states), $\DER_s \approx 1$ and no correction
      is needed; when the likelihood dominates (large states),
      $\DER_s$ can exceed~$3$ and the correction is substantial.
    \item \emph{Hyperparameters}
      $(\bGamma_k, \bSigma_\varepsilon)$\textup{:} no correction.
      The sandwich estimator is not applicable to variance components
      estimated from between-state variation; the DER concept does not
      extend to these parameters.
  \end{enumerate}
\end{remark}

%% ---------------------------------------------------------------------------
%%  T2-1: Approximation quality of block-wise approach
%% ---------------------------------------------------------------------------

\begin{remark}[Approximation quality of the block-wise correction]
\label{smc:rem-blockwise-quality}
  The block-wise application ignores the cross-block covariance between
  fixed effects $\boldsymbol{\xi}$ and random effects $\bepsilon_s$ in
  $\mathbf{V}_{\mathrm{sand}}$.  The approximation is essentially exact
  when the prior dominates ($\ESS_s \ll q$), since $\bepsilon_s$ is then
  nearly independent of the data.  When data dominate ($\ESS_s \gg q$),
  the cross-block covariance can be non-negligible; the sensitivity
  comparison in~\cref{app:extended-results} addresses this case.
\end{remark}

%% ---------------------------------------------------------------------------
%%  T1-4: MCMC vs design-consistent shrinkage
%% ---------------------------------------------------------------------------

\begin{remark}[MCMC versus design-consistent shrinkage]
\label{smc:rem-mcmc-vs-design}
  The shrinkage factor
  $B_s^{(w)}
  = \bSigma_\varepsilon(\bSigma_\varepsilon + \mathbf{V}_s^{-1})^{-1}$
  describes the \emph{design-consistent} (post-sandwich-correction)
  posterior behavior, not the raw MCMC pseudo-posterior.  The MCMC
  sampler sees precision proportional to
  $W_s = \sum_{i \in s}\tilde{w}_i$, which exceeds the effective
  sample size $\ESS_s = W_s^2 / \sum_{i \in s}\tilde{w}_i^2$ whenever
  within-state weights are heterogeneous.  The Cholesky affine
  transformation~\eqref{eq:cholesky-transform} maps the MCMC output to
  the design-consistent target, inflating the variance by the factor
  $\DER_s$ (\cref{smc:der}).  The gap $W_s - \ESS_s$ quantifies the
  over-concentration of the raw pseudo-posterior relative to the
  design-aware target\textup{:} large gaps (common in states with a
  few heavily weighted providers) signal that the naive MCMC credible
  intervals are too narrow and that the sandwich correction is most
  important.
\end{remark}

%% ---------------------------------------------------------------------------
%%  Limitations
%% ---------------------------------------------------------------------------

\begin{remark}[Non-Gaussian posterior shapes]
\label{smc:rem-nongaussian}
  The Cholesky transformation is an affine operation that preserves
  only first and second moments.  If the pseudo-posterior is
  substantially non-Gaussian (skewed or multimodal), the transformed
  draws will have the correct covariance but may misrepresent the
  posterior shape.  For the HBB model, the fixed effects
  $(\balpha,\bbeta)$ are expected to be approximately Gaussian by the
  Bernstein--von Mises result (\cref{smc:thm-bvm}) given
  $N = 6{,}809$; the log-dispersion $\log\kappa$ is also
  well-approximated by a normal.  For variance components
  ($\bSigma_\varepsilon$), the posterior may exhibit right-skewness,
  but these parameters are not subject to the Cholesky correction
  (Block~3 above).
\end{remark}

\begin{remark}[Positive definiteness requirements]
\label{smc:rem-posdef}
  The transformation requires both $\bSigma_{\mathrm{MCMC}}$ and
  $\mathbf{V}_{\mathrm{sand}}$ to be positive definite so that their
  Cholesky factors exist.  For $\bSigma_{\mathrm{MCMC}}$, positive
  definiteness holds almost surely when the number of MCMC draws
  satisfies $M \gg \dim(\boldsymbol{\xi}) = 11$; in our application
  $M = 4{,}000$ post-warmup draws.  For
  $\mathbf{V}_{\mathrm{sand}}$, positive definiteness requires the
  effective degrees of freedom to exceed the parameter dimension:
  $\sum_{h=1}^{H}(C_h - 1) = 386 > 11$ (\cref{smc:rem-dof}), which
  is easily satisfied.
\end{remark}
